\subsection{二軸の全値判定と105点への必要被覆}
\label{subsec:cubic-third-root-unit-boundary-axes}
\begin{theorem}[軸上の全境界正常核]
\label{thm:cubic-third-root-unit-boundary-axes}
$E=0,D\ne0$ のH3候補は $D^3+\omega^2=0$ に限り、
$D=0,E\ne0$ のH3候補は $E^3+\omega=0$ に限る。
それぞれの例外点で、境界の中山固定正常三次空間は一次元であり、
原 $B/(x)$ 座標で次の元が生成する：
\[
 \Gamma_E=ycb-D^3byc+cby,\qquad
 \Gamma_D=ybc+bcy-E^3cyb.
\]
従って $w\ne0$ を正規化した全正常一次境界のH3必要候補は、
前節の99点と、この6軸点を合わせた105個の $(D,E)$ に限る。
\end{theorem}
\begin{proof}
原点はPHS条件で除外済みなので、一方の軸変数は非零としてよい。
各軸の三次元Sklyanin商の次数3,4を、元の全二次関係から構成する。
商次元は10,15で、全relation pivotは非零軸変数のmonomial単元である。
従ってこの計算は軸の全非零点に有効である。
中山固定三次基底4個に対する三正常式の行列は $11\times4$ となる。
行ごとに軸変数の冪を掛けて多項式行列に直すと、Smith因子はそれぞれ
\[
 (1,1,1,D^3+\omega^2),\qquad
 (1,1,1,E^3+\omega).
\]
左右変換の行列式は非零定数であり、変換恒等式も元行列へ掛け戻して確認する。
従って記載した例外以外では境界正常核は零で、全正常三次は $x$ 因子を持ちH3でない。
例外ではrank3、核は一本である。

六次体 $k_0=\mathbb Q[t]/(t^6-t^3+1)$ で新しく元関係を再構成し、
$E=0$ では $D=t,\omega=t^3-1$、$D=0$ では $E=t,\omega=-t^3$ とする。
全六埋込みで元正常式を再代入して、表示した生成元と核一次元を確認する。
再現器は \texttt{verify\_cubic\_third\_root\_unit\_boundary\_axes.py}。
\end{proof}

\begin{proposition}[六軸点の全変形正常三次]
\label{prop:cubic-third-root-unit-boundary-axis-incidence}
六軸点では $D+\omega E\ne0$ なので、既知のshearで $Z=0$ とできる。
CY untwist $C$ の座標で、H3の候補正常三次は全 $u,v$ について次の一本に限る：
\begin{align*}
 E=0:\quad G_E={}&ycb+\omega^2byc+cby
           +v(\omega^2bc+cb)x-uD^2yx^2,\\
 D=0:\quad G_D={}&ybc+bcy+\omega cyb
           +v(bc+\omega cb)x-uE^2yx^2.
\end{align*}
ここで正常化は $\delta=(1,1,\omega^2,\omega)$ である。
この正常核の認識だけから、商の整域性やH3許容性は結論しない。
\end{proposition}
\begin{proof}
上の六次体を係数体とし、$u,v$ を両方残す。
元の二次関係をuntwistし、その次数3,4の全商を多項式係数で構成する。
全relation pivotはパラメータ定数であり、全特殊化で同じ語基底を使える。
中山固定三次8係数の正常行列は各 $26\times8$。
パラメータ定数単元だけによる7回の消去で全残余が零となり、核は自由階数1となる。
表示元の自由座標は非零定数なので、全特殊化後もその元が核全体を生成する。
さらに元の四生成元の正常式へ再代入して、多項式として零を確認する。
再現器は \texttt{verify\_cubic\_third\_root\_unit\_boundary\_}\\
\texttt{axis\_incidence.py}。
\end{proof}

\begin{corollary}[軸切断の残る最小パラメータ]
$Z=0,u=v=0$ の六軸切断はH3でない。
残る軸切断は $v=1,u=a$ の一変数族と、$v=0,u=1$ の固定切断を調べればよい。
\end{corollary}
\begin{proof}
$u=v=0$ では唯一の正常三次は純baseにあり、包絡はその三次元基底上の正常一次Ore拡大である。
従って補題\ref{lem:cubic-skew-cone-vertex}の頂点障害でH3でない。
$x$ のscalar変換は $(u,v)$ に重み $(2,1)$ で作用する。
$v\ne0$ なら $v=1,u=a$ へ、$v=0,u\ne0$ なら $u=1$ へ移せる。
この操作は表示した正常三次も同時に移す。
\end{proof}

\begin{theorem}[非軸99点の排除と六点への必要被覆]
\label{thm:cubic-third-root-unit-boundary-six-points}
前節の非軸99点では、中山固定境界正常三次の核は全て零である。
従って $w\ne0$ の正常一次chartのH3候補は、固定した $\omega$ ごとに
\[
 (E=0,\ D^3=-\omega^2),\qquad (D=0,\ E^3=-\omega)
\]
の六点に限る。これは $(D,E)$ の必要被覆であり、各点の全 $u,v$ 切断が
H3を与えるという主張ではない。
\end{theorem}
\begin{proof}
99点を与える11次式は、二つの一次因子と三つの三次因子の積である。
一次因子は $\mathbb Q(\omega)$ 上で扱う。
三次因子は $K=\mathbb Q[q]/(q^6+q^3+1)$ 上の次の三表示で扱う：
\[
\begin{array}{c|c|c}
 \omega&t&\text{消滅する三次因子}\\\hline
 q^6&(q-1)^{-1}&t^3+(1-\omega)(t^2+t+1/3)\\
 q^3&(q-1)^{-1}&t^3+(2+\omega)(t^2+t+1/3)\\
 q^3&q+q^{-1}-1&t^3+3t^2-1.
\end{array}
\]
$q^6+q^3+1$ は既約である。最初の二表示では $q=1+t^{-1}$、
最後の表示では $t$ の最小多項式は記載した有理既約三次式である。
従って固定した $\omega$ を保つ三埋込みは、それぞれの三次因子の全根を覆う。
これは体の埋込みによる全点の検算であり、幾何学的同型による同一視ではない。

各表示と $a,b=0,1,2$ に対し
$D=-\omega^b(t+1), E=\omega^{b-a}t$ を代入する。
元の三つの二次関係から次数3,4商を作ると、次元はそれぞれ10,15となる。
中山固定三次基底は4個で、その三正常式の行列の4列は全て独立である。
公開再現器は45ケース全てで非零 $4\times4$ minorを再計算し、
その行番号と行列式の体基底係数を証明書に記録する。
各一次ケースは1点、三次ケースは3点を覆い、点数は $9(1+1+3+3+3)=99$。
従って全点で境界正常核が零である。
H3正常三次の境界像が零なら $x$ 因子を持ち、整域条件に反する。
再現器は \texttt{verify\_cubic\_third\_root\_unit\_boundary\_99.py}。
軸の定理と合わせて結論を得る。
\end{proof}

\begin{theorem}[六軸点の全切断排除]
\label{thm:cubic-third-root-unit-boundary-axes-no-H3}
上の六軸点は、全 $Z,u,v$ に対してH3を与えない。
従って六係数の正常一次chartの $w\ne0$ 部分は全てH3から除外される。
\end{theorem}
\begin{proof}
まず $E=0,D^3=-\omega^2$ とし、既知のshearにより $Z=0$ とする。
H3候補と仮定し、$A=C/(G_E)$ と置く。すると
$H_A=(1-t^3)/(1-t)^4$、$x$ は中央正則である。
次の中央三次を用いる：
\[
 \Psi=y^3+vxy^2+ux^2y.
\]
全 $u,v$ を残した多項式係数で、$A$ の次数9,10を全元関係から構成する。
中山重み零の三次ではなく九次の語基底46個について、四交換子の行列を作ると、
パラメータ定数単元だけの41回の消去で全残余が零となる。
自由階数5の核の全生成元を元交換子へ再代入する。
そのうち、$x=0$ で $y^9$ から独立な元 $z_9$ を選べる。
$z_9\bmod x$ の係数は全て $u,v$ に依存しない。
$\Psi$ の四交換子も元関係で零である。

公開再現器と証明書に記載した $z_9$ に対して、
$A/(x,\Psi,z_9)$ の全関係からの次数0から13の商次元は
\[
 1,3,6,8,9,9,9,9,9,8,6,3,1,0.
\]
次数13で零なので以後も零である。この商は $u,v$ に依存しない。
全元関係の左右作用も再挿入して確認する。
従って $A$ は三中心元の部分環上有限生成である。
$\operatorname{GKdim}A=3$ なので三中心元は代数的独立である。
商の語基底の持上げによる自由加群の全射は、
上のHilbert級数と
$(1-t)(1-t^3)(1-t^9)H_A$
との一致から同型である。
従って $A$ は $S_0=k[x,\Psi,z_9]$ 上自由階数81となる。
この議論はH3から必要となるHilbert級数を使っており、
有限次数計算だけから包絡の全次数AS性を推論していない。

元の $x=0$ 上には、全パラメータに共通の中心fiber
\[
 F=A/(x,\Psi-1,z_9)
\]
がある。これが次数9の中心単純代数であることを、正の有限体witnessで確認する。
係数体 $\mathbb Q[t]/(t^6-t^3+1)$ の $t=4\pmod {109}$ では
$\omega=63$ となる。元の $x$ とは別の中央変数 $h$ でこのfiberを斉次化し、
$y^3=h^3,z_9=0$ を課す。
次数12,13の次元はともに81、$h$ の作用は可逆である。
$h=1$ の正則表示で、中心交換子行列にrank80の非零minorがあり、
正則trace Gram行列の行列式は $108\pmod {109}$ で非零となる。
全元関係の81次行列への再代入も零である。
計算は浮動小数点を使わず、各積の整数上界は $81\cdot108^2<2^{63}$。

自由階数81が標数0で既知なので、上の正の商基底minorは標数0でも同じ基底を与え、
中心minorとtrace行列式も非零である。
従って標数0の $F$ は中心単純、次数9である。
体の埋込みによって六つの全埋込みに同じ結論が成り立つ。
このfiberは全 $u,v$ で同一なので、全パラメータで中心一般点がCSA次数9となる。
自由性と一般点単純性から $A$ は素、整閉性から $Z(A)=S_0$。
$U=(A[x^{-1}])_0$ は素、全中心は $k[\psi,z]$、PI次数9となる。

正常三次商の双対公式と中心自由性から
\[
 \operatorname{Hom}_{S_0}(A,S_0)=A_\delta(12),\qquad
 \operatorname{Hom}_{k[\psi,z]}(U,k[\psi,z])=U_\delta,\quad\delta^3=1.
\]
従って補題\ref{lem:cubic-Frobenius-order-hereditary}と
命題\ref{prop:cubic-Frobenius-order-period}の証明がそのまま適用できる。
H3なら有限大域次元から周期は3を割り、整域性と次数9から曲面period--index定理により
周期9となる。これは矛盾である。

最後に原表示を反対環にすると、生成元を変えずに
\[
 B(D,E,0,u,v;\omega)^{\mathrm{op}}
   =B(E,D,0,u,v;\omega^2)
\]
となる。各二次関係の語順を逆転すればこの等式を直接確認できる。
しかも $G_E^{\mathrm{op}}$ は右辺の $G_D$ に一致する。
整域性と左右有限大域次元は反対環で保存されるため、
既に証明した $E=0$ の矛盾は $D=0$ にも適用できる。
これは反対環をintrinsic moduliの同型として同一視する議論ではない。

再現器は \texttt{verify\_cubic\_third\_root\_unit\_boundary\_}\\
\texttt{axis\_center.py} と同接頭辞の \texttt{axis\_CSA.py}。
proper楕円上の被約九点軌道の議論を、この特異境界へ仮定なしに外挿してはいない。
\end{proof}
